\begin{problem}[The continuous exponential short exact sequence]\label{prob:vi16-p04}
On a topological space $X$, analyze the morphism $\exp(2\pi i-):\mathcal C^0_{\mathbb R}\to\mathcal C^0_{S^1}$.  Prove the resulting sequence with kernel $\underline{\mathbb Z}$ is exact as sheaves and explain the failure of global surjectivity on $S^1$.
\end{problem}

\ifdefined\FullProblemDossiers
\paragraph{Why this problem matters.}
This is the cleanest example separating local exactness from global lifting.

\paragraph{Useful ideas.}
Use local arguments of circle-valued maps and the degree of the identity map.

\paragraph{Guided hints.}
\begin{enumerate}
\item Identify the kernel pointwise.
\item Use a small arc in $S^1$ to choose a continuous argument.
\item For global failure, use winding number.
\end{enumerate}

\begin{solution}
The kernel consists of continuous integer-valued functions, hence locally constant $\mathbb Z$-valued functions.  Every circle-valued germ admits a continuous argument on a sufficiently small neighborhood, so the exponential stalk map is surjective.  The stalkwise criterion gives exactness.  On $X=S^1$, the identity circle map has degree one and therefore cannot be the exponential of a global continuous real-valued function, whose exponential has degree zero.
\end{solution}

\paragraph{What happened mathematically.}
The quotient sheaf sees local logarithms and forgets the winding obstruction that prevents a global logarithm.

\paragraph{Extensions.}
Replace continuous by holomorphic functions to obtain P5.

\paragraph{Related problems.}
VI.16.P6 and P7 use this sequence to show failure of right exactness for global sections/direct image.

\paragraph{Provenance.}
New synthesis from preferred corpus examples; no separate numbered legacy problem.
\else
\paragraph{Provenance.} New synthesis from preferred corpus examples; no separate numbered legacy problem.
\fi
